\documentclass[12pt]{article}

\usepackage[margin=1in]{geometry}
\usepackage{amsmath, amssymb}
\usepackage{enumerate}
\usepackage{hyperref}

\begin{filecontents*}[overwrite]{refs.bib}
@inproceedings{lin2017,
  author    = {Tsung-Yi Lin and Priya Goyal and Ross Girshick and Kaiming He and Piotr Doll{\'a}r},
  title     = {Focal Loss for Dense Object Detection},
  booktitle = {ICCV},
  year      = {2017},
}
@inproceedings{cui2019,
  author    = {Yin Cui and Menglin Jia and Tsung-Yi Lin and Yang Song and Serge Belongie},
  title     = {Class-Balanced Loss Based on Effective Number of Samples},
  booktitle = {CVPR},
  year      = {2019},
}
@inproceedings{cao2019,
  author    = {Kaidi Cao and Colin Wei and Adrien Gaidon and Nikos Ar{\'{e}}chiga and Tengyu Ma},
  title     = {Learning Imbalanced Datasets with Label-Distribution-Aware Margin Loss},
  booktitle = {NeurIPS},
  year      = {2019},
}
@article{wang2019,
  author    = {Shuo Wang and Wanli Liu and Lingqiao Wu and Lihua Cao and Qi Meng and Wei Chen},
  title     = {Learning from Imbalanced Data Sets with Weighted Cross-Entropy Function},
  journal   = {Pattern Recognition Letters},
  year      = {2019},
}
@inproceedings{he2016,
  author    = {Kaiming He and Xiangyu Zhang and Shaoqing Ren and Jian Sun},
  title     = {Deep Residual Learning for Image Recognition},
  booktitle = {CVPR},
  year      = {2016},
}
@article{chawla2002,
  author    = {Nitesh V. Chawla and Kevin W. Bowyer and Lawrence O. Hall and W. Philip Kegelmeyer},
  title     = {{SMOTE}: Synthetic Minority Over-sampling Technique},
  journal   = {Journal of Artificial Intelligence Research},
  volume    = {16},
  pages     = {321--357},
  year      = {2002},
}
@techreport{krizhevsky2009,
  author      = {Alex Krizhevsky},
  title       = {Learning Multiple Layers of Features from Tiny Images},
  institution = {University of Toronto},
  year        = {2009},
}
@inproceedings{liu2019,
  author    = {Ziwei Liu and Zhongqi Miao and Xiaohang Zhan and Jiayun Wang and Boqing Gong and Stella X. Yu},
  title     = {Large-Scale Long-Tailed Recognition in an Open World},
  booktitle = {CVPR},
  year      = {2019},
}
@inproceedings{van_horn2018,
  author    = {Grant Van Horn and Oisin Mac Aodha and Yang Song and Yin Cui and Chen Sun and Alex Shepard and Hartwig Adam and Pietro Perona and Serge Belongie},
  title     = {The {iNaturalist} Species Classification and Detection Dataset},
  booktitle = {CVPR},
  year      = {2018},
}
@article{he2009,
  author    = {Haibo He and Edwardo A. Garcia},
  title     = {Learning from Imbalanced Data},
  journal   = {IEEE Transactions on Knowledge and Data Engineering},
  volume    = {21},
  number    = {9},
  pages     = {1263--1284},
  year      = {2009},
}
@article{buda2018,
  author    = {Mateusz Buda and Atsuto Maki and Maciej A. Mazurowski},
  title     = {A Systematic Study of the Class Imbalance Problem in Convolutional Neural Networks},
  journal   = {Neural Networks},
  volume    = {106},
  pages     = {249--259},
  year      = {2018},
}
@article{johnson2019,
  author    = {Justin M. Johnson and Taghi M. Khoshgoftaar},
  title     = {Survey on Deep Learning with Class Imbalance},
  journal   = {Journal of Big Data},
  volume    = {6},
  number    = {1},
  pages     = {27},
  year      = {2019},
}
@inproceedings{drummond2003,
  author    = {Chris Drummond and Robert C. Holte},
  title     = {{C4.5}, Class Imbalance, and Cost Sensitivity: Why Under-Sampling Beats Over-Sampling},
  booktitle = {Workshop on Learning from Imbalanced Datasets, ICML},
  year      = {2003},
}
@inproceedings{ren2018,
  author    = {Mengye Ren and Wenyuan Zeng and Bin Yang and Raquel Urtasun},
  title     = {Learning to Reweight Examples for Robust Deep Learning},
  booktitle = {ICML},
  year      = {2018},
}
@inproceedings{kang2020,
  author    = {Bingyi Kang and Saining Xie and Marcus Rohrbach and Zhicheng Yan and Albert Gordo and Jiashi Feng and Yannis Kalantidis},
  title     = {Decoupling Representation and Classifier for Long-Tailed Recognition},
  booktitle = {ICLR},
  year      = {2020},
}
@inproceedings{zhou2020,
  author    = {Boyan Zhou and Quan Cui and Xiu-Shen Wei and Zhao-Min Chen},
  title     = {{BBN}: Bilateral-Branch Network with Cumulative Learning for Long-Tailed Visual Recognition},
  booktitle = {CVPR},
  year      = {2020},
}
@inproceedings{menon2021,
  author    = {Aditya Krishna Menon and Sadeep Jayasumana and Ankit Singh Rawat and Himanshu Jain and Andreas Veit and Sanjiv Kumar},
  title     = {Long-Tail Learning via Logistic Regression},
  booktitle = {ICML},
  year      = {2021},
}
@inproceedings{tan2020,
  author    = {Jingru Tan and Changbao Wang and Buyu Li and Quanquan Li and Wanli Ouyang and Changqing Yan and Junjie Yan},
  title     = {Equalization Loss for Long-Tailed Object Recognition},
  booktitle = {CVPR},
  year      = {2020},
}
@inproceedings{zhang2021,
  author    = {Songyang Zhang and Zeming Li and Shipeng Yan and Xuming He and Jian Sun},
  title     = {Distribution Alignment: A Unified Framework for Long-Tail Visual Recognition},
  booktitle = {CVPR},
  year      = {2021},
}
@inproceedings{akiba2019,
  author    = {Takuya Akiba and Shotaro Sano and Toshihiko Yanase and Takeru Ohta and Masanori Koyama},
  title     = {Optuna: A Next-Generation Hyperparameter Optimization Framework},
  booktitle = {KDD},
  year      = {2019},
}
\end{filecontents*}

\usepackage[numbers]{natbib}

\title{Log-Scale Class Weighting for Robust Imbalanced Classification}
\author{Seung Ho Go \and Hyunseok Lee \and Jin Hee Yoon\\
        \small Sejong University}
\date{}

\begin{document}

\maketitle

\begin{abstract}
Class imbalance is a pervasive challenge in machine learning, arising across
diverse domains such as medical diagnosis, fraud detection, and visual recognition,
where minority classes are critically important yet severely underrepresented.
Weighted Cross-Entropy (WCE), the most direct remedy, assigns inverse-frequency
weights to each class, but these weights diverge rapidly as the imbalance ratio
grows, leading to numerical instability.
More sophisticated alternatives within the re-weighting family, such as
Class-Balanced (CB) Loss and Focal Loss, introduce additional hyperparameters
and can degrade below standard cross-entropy (CE) performance at high imbalance
ratios.
We propose \textbf{Logarithmic Weighted Cross-Entropy (LWCE)}, which replaces
inverse-frequency weights with logarithmically compressed counterparts
$w_i = 1/\log(n_i+1)$, and its generalization \textbf{PLWCE}
($w_i = 1/(\log(n_i+1))^\alpha$), which allows dataset-adaptive tuning of the
focusing strength via a single exponent $\alpha$ selected by Optuna grid search
on a small proxy subset.
We conduct a controlled evaluation on CIFAR-10-LT and CIFAR-100-LT at imbalance
ratios of 10, 50, and 100, using ResNet-32 with five random seeds.
Comparing six loss functions (CE, PWCE, LWCE, PLWCE, CB Loss, Focal Loss),
we find that LWCE and PLWCE consistently outperform CE on macro-averaged F1 and
tail-class accuracy as the imbalance ratio increases, while CB Loss and Focal Loss
degrade relative to CE at IR\,=\,100.
On CIFAR-10-LT at IR\,=\,100, PLWCE achieves F1-Macro of $0.458\pm0.018$
versus $0.431\pm0.021$ for CE and $0.406\pm0.015$ for CB Loss.
On CIFAR-100-LT at IR\,=\,100, LWCE and PLWCE reach F1-Macro of 0.185 and 0.186,
surpassing CE (0.179) and CB Loss (0.158).
These results suggest that logarithmic weight compression offers a simple, stable,
and parameter-efficient alternative to existing re-weighting strategies for
class-imbalanced classification.
\end{abstract}

\section{Introduction}
\label{sec:intro}

In many real-world classification tasks---spanning medical diagnosis, financial
fraud detection, network intrusion detection, and visual recognition---sample
counts across categories follow a long-tailed distribution: a small number of
majority classes concentrate the bulk of training data while many minority classes
are represented by only a handful of examples~\cite{cui2019,cao2019}.
Under such conditions, a classifier trained with standard cross-entropy (CE) loss is
implicitly biased toward majority classes: the aggregate gradient signal is dominated
by numerically abundant samples, so the model achieves high overall accuracy by
largely ignoring rare but equally valid categories---a phenomenon commonly termed
\textit{majority bias}.

Existing strategies to combat class imbalance fall broadly into two families.
\textbf{Resampling methods} directly manipulate the training distribution:
oversampling techniques such as SMOTE~\cite{chawla2002} synthesize new minority-class
instances, while undersampling discards majority-class examples to restore balance.
Although conceptually straightforward, both directions carry inherent risks.
Oversampling can introduce synthetic artifacts and lead to overfitting on
interpolated points, whereas undersampling discards potentially informative
majority-class samples, distorting the true data distribution.
Algorithm-level approaches avoid data manipulation altogether by modifying the
training objective to reflect class importance.
These include margin-based methods that widen decision boundaries for minority
classes~\cite{cao2019}, logit-adjustment techniques that calibrate class-prior
probabilities at inference time, and \textit{loss re-weighting} methods that
scale the cross-entropy loss by class-dependent weights.
While the former two families require architectural modifications or post-hoc
calibration, re-weighting methods are architecture-agnostic and applicable across
data modalities with minimal implementation overhead.
This paper focuses on the re-weighting family.

\textbf{Weighted Cross-Entropy (WCE)}~\cite{wang2019} scales the loss for class
$i$ by $w_i \propto 1/n_i$, where $n_i$ is the number of training samples in
that class.
Although conceptually simple, WCE suffers from a critical numerical instability:
as the imbalance ratio (IR) grows, the inverse-frequency weight $1/n_i$ diverges
rapidly, producing extreme gradient magnitudes that destabilize optimization.
This becomes especially problematic at severe imbalance---for instance, in
CIFAR-100-LT at IR\,=\,100 the rarest class contains only $n_i = 5$ samples,
making straightforward WCE impractical.

Several more sophisticated alternatives have been proposed.
\textbf{Focal Loss}~\cite{lin2017} addresses the class imbalance problem at the
\textit{sample} level by introducing a modulating factor $(1-p_t)^\gamma$ that
down-weights easy, often majority-class examples.
However, its behavior depends on per-sample prediction probabilities rather than
class frequency, making it sensitive to the calibration quality of the model and
yielding inconsistent gains across different imbalance regimes.
\textbf{Class-Balanced (CB) Loss}~\cite{cui2019} grounds re-weighting in the
concept of \textit{effective number of samples}, providing a theoretically
motivated schedule through a hyperparameter $\beta \in (0,1)$.
Yet $\beta$ requires careful tuning, and in our experiments CB Loss consistently
underperforms CE at high imbalance ratios---degrading to F1-Macro of 0.406 against
CE's 0.431 on CIFAR-10-LT at IR = 100, and 0.158 against 0.179 on CIFAR-100-LT
at the same ratio---suggesting that its weight schedule becomes too aggressive when
tail-class sample counts are very small.

To overcome the weight explosion of WCE while retaining its simplicity and
interpretability, we propose \textbf{Logarithmic Weighted Cross-Entropy (LWCE)},
which replaces the inverse-frequency weight with a logarithmically compressed
counterpart:
\begin{equation}
  w_i = \frac{1}{\log(n_i + 1)}.
  \label{eq:lwce}
\end{equation}
The logarithmic compression fundamentally bounds weight growth:
whereas WCE weights grow as $\mathcal{O}(n_i^{-1})$, LWCE weights grow only as
$\mathcal{O}((\log n_i)^{-1})$.
The method preserves a meaningful boost for rare classes without
inducing numerical instability, and requires no additional hyperparameters.
We further generalize LWCE to \textbf{Power-Log Weighted CE (PLWCE)} by introducing
a single tunable exponent $\alpha$:
\begin{equation}
  w_i = \frac{1}{\bigl(\log(n_i + 1)\bigr)^\alpha},
  \label{eq:plwce}
\end{equation}
which spans a continuous spectrum from near-uniform weighting ($\alpha \to 0$) to
aggressive minority focus ($\alpha \gg 1$), with LWCE recovered at $\alpha = 1$.
The optimal $\alpha$ per imbalance ratio is identified via Optuna grid search on a
20\% proxy subset, requiring negligible additional computation.

We evaluate the proposed methods on the standard long-tailed image classification
benchmarks, CIFAR-10-LT and CIFAR-100-LT, across three imbalance ratios
(IR $\in \{10, 50, 100\}$) using ResNet-32 trained from scratch~\cite{he2016}.
All experiments are repeated with five random seeds to report mean $\pm$ standard
deviation.
We compare LWCE and PLWCE against CE, Power-Scaled WCE (PWCE), CB Loss, and Focal
Loss using overall Top-1 accuracy, macro-averaged F1-score, and group-wise accuracy
on head (Many), middle (Medium), and tail (Few) classes defined by a tertile split
of training counts.

Results demonstrate that LWCE and PLWCE consistently improve over CE on F1-Macro
and Few-class accuracy as the imbalance ratio increases, while avoiding the
performance degradation exhibited by CB Loss and Focal Loss at high IRs.
On CIFAR-10-LT at IR = 100, PLWCE achieves F1-Macro of $0.458 \pm 0.018$,
surpassing CE ($0.431 \pm 0.021$) and CB Loss ($0.406 \pm 0.015$).
On CIFAR-100-LT at IR = 100, LWCE and PLWCE reach F1-Macro of $0.185$ and $0.186$,
outperforming CE ($0.179$) and CB Loss ($0.158$), with the gap widening at the
tail-class level.

The main contributions of this paper are:
\begin{enumerate}
  \item We formally characterize the \textit{weight explosion} failure mode of WCE
        and demonstrate its practical consequences on CIFAR-LT benchmarks.
  \item We propose LWCE, a parameter-free drop-in replacement for WCE based on
        logarithmic weight compression, and its single-parameter generalization PLWCE.
  \item We provide a controlled comparison of six loss functions on CIFAR-10-LT and
        CIFAR-100-LT with five-seed repetition, revealing that logarithmic re-weighting
        offers a more stable and consistent improvement over CE than CB Loss or Focal
        Loss, particularly in the high-imbalance regime.
\end{enumerate}

The remainder of this paper is organized as follows.
Section~\ref{sec:related} reviews related work on long-tailed learning.
Section~\ref{sec:method} formally defines LWCE and PLWCE and analyzes their
theoretical properties.
Section~\ref{sec:experiments} describes the experimental setup and presents results.
Section~\ref{sec:conclusion} concludes.

\section{Related Work}
\label{sec:related}

\section{Proposed Method}
\label{sec:method}

\section{Experiments}
\label{sec:experiments}

\section{Conclusion}
\label{sec:conclusion}

\nocite{*}
\bibliographystyle{plainnat}
\bibliography{refs}

\end{document}
